% AC-S51: actual parity-separated spectral-edge comparison.
\section{An actual ground-line inverse from the parity blocks}
\label{sec:spectral-parity-edge}

The comparison here is between two specified operators on the actual
localized Weil spaces.  It does not assert a lower bound for either
parity minimum.  Its eventual application to the original \(\Xi\)
retains zero simplicity explicitly.  Without that assumption the
limit obtained below is the product over distinct ordinates.

Assume \(A\).  Use
\[
 \widehat f(t)=\int_{\mathbb R}f(x)e^{-itx}\,dx,\qquad
 D=-i\partial_x,\qquad
 q(f,g)=\sum_{\gamma\in\Gamma}
       m_\gamma\widehat f(\gamma)\overline{\widehat g(\gamma)} .
 \tag{PE.1}
\]
All inner products are linear in the first variable.  Let \(H\)
be the global Weil completion, and let \(H_L\) be the local
completion of \(C_c^\infty(-L,L)\), embedded isometrically in \(H\).
The local strict positivity and compact resolvent of the actual Weil
operator \(\mathscr A_L\), as in AC-S03, identify \(H_L\) with its
form domain and give a compact inclusion
\[
 J_L:H_L\longrightarrow L^2(-L,L).
\]
Write \(\Pi_L\) for the Weil-orthogonal projection onto \(H_L\).

For completeness the compact smooth core can be checked directly
in the same actual form domain.  The fixed-window explicit formula
makes its form norm equivalent to the norm with Fourier weight
\(1+\log(2+|t|)\); the prime and polar terms are bounded on the
window.  For a zero-extended form-domain function, first shrink
support by the dilation \(f_r(x)=r^{-1/2}f(x/r)\), \(r<1\).
Dilation is strongly continuous for this weighted Fourier norm:
its operators are uniformly bounded for \(1/2\le r\le1\), and
continuity first holds on Schwartz functions and then by density.
Mollification at a scale smaller than the remaining boundary
margin now gives compact smooth approximants in the same norm.
This proves the required core assertion, with constants allowed
to depend on the fixed window.

The elementary sampling-surjectivity proof of AC-S50 identifies
\(H\) with \(\ell^2(\Gamma,m)\).  Put \(C=D^{-1}\), meaning the
selfadjoint diagonal operator with entries \(\gamma^{-1}\).
It is Hilbert--Schmidt, and
\[
 C_L=\Pi_LC\Pi_L
 \tag{PE.2}
\]
is selfadjoint on \(H_L\); when used on \(H\), it is extended by zero
on \(H_L^\perp\).  Reflection gives the orthogonal decompositions
\(H=H^+\oplus H^-\) and \(H_L=H_L^+\oplus H_L^-\).
In the latter decomposition write
\[
 C_L=
 \begin{pmatrix}0&X_L^*\\X_L&0\end{pmatrix},
 \qquad
 T_L=J_L^*J_L=
 \begin{pmatrix}T_{L,+}&0\\0&T_{L,-}\end{pmatrix}.
 \tag{PE.3}
\]
The adjoints in (PE.2)--(PE.3) are Weil adjoints, not spatial
restriction or ordinary \(L^2\) adjoints.

Let \(a_e(L)\) and \(a_o(L)\) be the actual even and odd minima of
\(\mathscr A_L\).  Both are positive.  The whole-window minimum is
their minimum.  We will work at windows satisfying
\[
                 a_e(L)<a_o(L).
 \tag{PE.4}
\]
Put \(\alpha=a_e(L)\), \(M_L=\alpha T_L\), and
\[
 S_+=I-M_{L,+},\qquad S_-=I-M_{L,-},\qquad
 \kappa_L=1-\frac{a_e(L)}{a_o(L)}>0 .
 \tag{PE.5}
\]
The operator \(S_-\) is bounded below by \(\kappa_L I\).
The kernel of \(S_+\) is the even bottom eigenspace.

\begin{theorem}[AC-S51: parity-separated actual edge comparison]
\label{AC-S51}
At every window satisfying (PE.4), the actual minimum is simple
and even.  If \(e_L\) is its Weil-unit eigenvector, then
\(\ell_L(e_L):=\int_{-L}^L e_L\ne0\).
On the Hilbert quotient of \(H_L\) by \(\mathbb C e_L\), with form
\(q-\alpha\langle\cdot,\cdot\rangle_2\), there is an injective compact
selfadjoint operator \(\mathcal K_L\) such that
\[
 \mathcal K_L[Df]=[f]\qquad(f\in C_c^\infty(-L,L)).
 \tag{PE.6}
\]
Its nonzero inverse eigenvalues are precisely the zeros of
\(\widehat e_L\).  Those zeros are real and simple.

On the original Weil space this operator is unitarily represented,
with an additional zero on \(\mathbb C e_L\), by
\[
 B_L=
 \begin{pmatrix}
  0&S_+^{1/2}X_L^*S_-^{-1/2}\\
  S_-^{-1/2}X_LS_+^{1/2}&0
 \end{pmatrix}.
 \tag{PE.7}
\]
In particular no inverse of the singular even block is used in this
formula.

Suppose (PE.4) holds for all sufficiently large \(L\), and suppose
there is a fixed \(\kappa>0\) such that
\[
       \kappa_L\ge\kappa
       \quad\hbox{for all sufficiently large }L.
 \tag{PE.8}
\]
Then, after extension by zero to \(H\),
\[
 \|B_L-C\|_{\mathcal S_2}\longrightarrow0,\qquad
 \frac{\widehat e_L(z)}{\widehat e_L(0)}
       \longrightarrow
       \Delta_\Gamma(z):=
          \prod_{\gamma>0,\ \gamma\in\Gamma}
                    (1-z^2/\gamma^2)
 \tag{PE.9}
\]
locally uniformly for \(z\in\mathbb C\).
The product in (PE.9) uses each distinct positive ordinate once.
If all nontrivial zeta zeros are simple, its value is
\(\Xi(z)/\Xi(0)\).
No \(L^2\) convergence of spatially unit-normalized ground vectors,
weighted \(L^1\) convergence, or effective convergence modulus is
asserted here.
\end{theorem}

\begin{proof}
\proofstep{1. The exact compact primitive on the mean-zero subspace.}
Let \(\ell=\ell_L\).  Local strict positivity makes \(\ell\)
a bounded functional on \(H_L\).
For a compact smooth mean-zero function \(f\), put
\[
       (\mathcal P f)(x)=i\int_{-L}^x f(y)\,dy .
 \tag{PE.10}
\]
It is compact smooth in the same interval, and
\(\widehat{\mathcal P f}(\gamma)=\widehat f(\gamma)/\gamma\).
Thus \(C_Lf=\mathcal P f\).  Also
\(C_LDg=g\) for every \(g\in C_c^\infty(-L,L)\).

The identity extends to every \(f\in H_L\) with \(\ell(f)=0\).
Indeed approximate \(f\) in the form norm by compact smooth
functions and remove their integral error using one fixed compact
smooth function of integral one inside the interval.
The corrected sequence converges both in \(H_L\) and in \(L^2\).
Its primitives converge in \(H^1_0(-L,L)\), while boundedness of
\(C_L\) gives convergence in \(H_L\).  The two limits have the same
\(L^2\) representative.  Therefore
\[
 f\in H_L,\ \ell(f)=0
 \quad\Longrightarrow\quad
 C_Lf=\mathcal P f\in H^1_0(-L,L),\quad DC_Lf=f .
 \tag{PE.11}
\]
In particular the primitive identity is a statement about the
actual local form domain, not a formal inverse on a completion
unrelated to compact support.

\proofstep{2. Simple even minimum and a nonzero integral.}
There is no nonzero mean-zero even bottom eigenvector.
Otherwise let \(u\) be such a vector and put \(v=C_Lu\).
By (PE.11), \(v\) is an odd compact primitive and is nonzero.
Its oddness implies \(\ell(v)=0\), so \(w=C_Lv\) is also a compact
primitive.  Selfadjointness of \(C_L\), the eigenvector form
identity, and integration by parts between \(v,w\in H^1_0\)
give
\[
 \begin{aligned}
 q(v,v)&=q(u,C_Lv)
      =\alpha\langle u,w\rangle_2
      =\alpha\langle Dv,w\rangle_2\\
      &=\alpha\langle v,Dw\rangle_2
       =\alpha\|v\|_2^2 .
 \end{aligned}
 \tag{PE.12}
\]
The minimum principle makes \(v\) an odd bottom vector,
contradicting (PE.4).
If the even bottom eigenspace had dimension at least two, a
nonzero linear combination would have integral zero, so this is
also impossible.  Thus the whole minimum is simple even and
\(\ell(e_L)\ne0\).

Write \(e=e_L\).  The form
\[
             q_\alpha(f,g)=q(f,g)-\alpha\langle f,g\rangle_2
 \tag{PE.13}
\]
has kernel exactly \(\mathbb C e\).  On its Weil-orthogonal
complement it is equivalent, for this fixed window, to \(q\),
because \(T_L\) is compact and its top eigenvalue is simple.
Consequently it defines a Hilbert norm on \(H_L/\mathbb C e\).
No lower bound uniform in \(L\) is used at this step.

\proofstep{3. The rank-one commutator and the quotient inverse.}
For \(f\in H_L^+\cap\ker\ell\) and \(g\in H_L^-\), both have
integral zero.  The compact primitive and ordinary integration
by parts give
\[
 \langle C_Lf,g\rangle_2=\langle f,C_Lg\rangle_2.
\]
Using the definition of \(T_L\), this says that
\[
       (X_LM_{L,+}-M_{L,-}X_L)f=0 .
\]
The displayed map therefore factors through \(\ell\) on the even
space.  Since \(M_{L,+}e=e\), its value at \(e\) is
\(S_-X_Le\).  As \(\ell(e)\ne0\), we obtain the exact identity
\[
 X_LM_{L,+}-M_{L,-}X_L
       =S_-X_Le\,\frac{\ell|_{H_L^+}}{\ell(e)}.
 \tag{PE.14}
\]
Equivalently,
\[
 X_L-X_Le\,\frac{\ell|_{H_L^+}}{\ell(e)}
                     =S_-^{-1}X_LS_+ .
 \tag{PE.15}
\]

Define on \(H_L\)
\[
 K_Lf=C_L\left(f-e\,\frac{\ell(f)}{\ell(e)}\right).
 \tag{PE.16}
\]
It is bounded and compact, kills \(e\), and hence induces an
operator \(\mathcal K_L\) on the quotient.  Formula (PE.11) says
that \(K_Lf\) is the actual compact primitive of the parentheses.
For compact smooth \(f\), \(\ell(Df)=0\), so (PE.6) follows.

Let \(Q_+\) be Weil-orthogonal projection off \(e\) in \(H_L^+\).
In the quotient coordinates
\(Q_+H_L^+\oplus H_L^-\), the induced operator has blocks
\(Q_+X_L^*\) and \(S_-^{-1}X_LS_+|_{Q_+H_L^+}\).
The map
\[
 [f]\longmapsto
   \bigl(S_+^{1/2}Q_+f_+,\,S_-^{1/2}f_-\bigr)
 \tag{PE.17}
\]
is an isometry from the \(q_\alpha\) quotient onto
\(Q_+H_L^+\oplus H_L^-\) with its \(q\) norm.
Its inverse is bounded at this fixed window.
Conjugation yields (PE.7), restricted off its zero vector \(e\).
The two displayed blocks are adjoints.  This proves compact
selfadjointness of \(\mathcal K_L\).

It remains to prove injectivity.
If \([K_Lf]=0\), then \(K_Lf=ce\).
If \(c\ne0\), (PE.16) gives \(e=C_Lg\) with
\(g=c^{-1}(f-e\ell(f)/\ell(e))\in H_L\cap\ker\ell\).
Injectivity of \(C_L\), and its reversal of parity, make \(g\)
odd.  Here injectivity follows at once from
\(C_LD\varphi=\varphi\) and density of compact smooth \(\varphi\).
By (PE.11), \(e=\mathcal P g\in H^1_0\) and \(De=g\).
For each \(\varphi\in C_c^\infty(-L,L)\),
\[
 \begin{aligned}
 q(g,\varphi)&=q(e,D\varphi)
     =\alpha\langle e,D\varphi\rangle_2\\
     &=\alpha\langle g,\varphi\rangle_2 .
 \end{aligned}
 \tag{PE.18}
\]
Both sides are continuous in the form norm of \(\varphi\);
core density extends the identity to every \(\varphi\in H_L\),
in particular to \(g\).  It would make the nonzero odd \(g\)
a bottom eigenvector, again a contradiction.
This argument does not assume that \(Dg\) belongs to the form
domain.  We conclude \(c=0\).
Then injectivity of \(C_L\) in (PE.16) implies
\(f\in\mathbb C e\), so \([f]=0\).

\proofstep{4. The spectrum is the actual ground Fourier zero set.}
Put \(F(z)=\widehat e(z)\).  It is even, entire of exponential
type at most \(L\), real on the real axis after choosing \(e\)
real, and \(F(0)=\ell(e)\ne0\).
For any fixed complex \(z\), Fourier evaluation is bounded on
\(H_L\), by its \(L^2\) inclusion and compact support.
If \(F(z)=0\), it annihilates \(\mathbb C e\), so descends to
a bounded nonzero functional on the \(q_\alpha\) quotient.

The primitive formula gives, for \(z\ne0\),
\[
 \widehat{K_Lf}(z)
        =\frac{\widehat f(z)-F(z)\ell(f)/\ell(e)}{z}.
 \tag{PE.19}
\]
At a zero of \(F\), evaluation therefore intertwines
\(\mathcal K_L\) with the scalar \(1/z\).
Its Riesz vector is a nonzero eigenvector with eigenvalue
\(1/\overline z\).  Selfadjointness shows \(z\) is real.

Conversely let \(\mathcal K_Lw=\beta w\) for a nonzero eigenvalue
\(\beta\).  The bounded functional on \(H_L\) defined by
\(f\mapsto q_\alpha([f],w)\) is a distribution on the interval.
Equation (PE.6) makes it satisfy
\(\Lambda(Df)=\beta^{-1}\Lambda(f)\) on compact smooth tests.
The elementary constant-coefficient distributional ODE gives
\(\Lambda(f)=c\,\widehat f(\beta^{-1})\).
The constant is nonzero.  Both functionals are continuous in
\(H_L\), and the left one annihilates \(e\); hence
\(F(\beta^{-1})=0\).
The same ODE shows every nonzero eigenspace is one-dimensional.

Finally \(F\) has no multiple zero.  If \(F(z)=F'(z)=0\),
both evaluation and its derivative descend to the quotient.
Differentiating (PE.19) gives, at this real nonzero \(z\),
\[
  \ell'_z(\mathcal K_L[f])
       =z^{-1}\ell'_z([f])-z^{-2}\ell_z([f]).
\]
For their Riesz vectors this is a generalized eigenvector relation
\[
       (\mathcal K_L-z^{-1})w'_z=-z^{-2}w_z .
 \tag{PE.20}
\]
Taking its inner product with the nonzero eigenvector \(w_z\)
contradicts selfadjointness.  Thus all zeros are simple.
No assertion about the multiplicities of the limiting zeta
zeros has entered this fixed-window argument.

\proofstep{5. The normalized metric tends strongly to zero.}
A compact smooth \(h\) with integral one, rescaled by
\(h_R(x)=R^{-1}h(x/R)\), satisfies
\[
 q(h_R)\le C R^{-2}\sum_\gamma m_\gamma/\gamma^2,
       \qquad \|h_R\|_2^2=R^{-1}\|h\|_2^2 .
\]
With \(R\) a fixed small multiple of \(L\), the minimum principle
gives \(a_1(L)\to0\).
At all windows under consideration \(\alpha=a_1(L)\), so the
extension of \(M_L=\alpha T_L\) to \(H\) is a positive contraction.
For every fixed compact smooth \(f\), at all sufficiently large
windows
\[
 \|M_Lf\|_q^2\le q(M_Lf,f)
                  =\alpha\|f\|_2^2\longrightarrow0 .
 \tag{PE.21}
\]
Density and contraction imply \(M_L\to0\) strongly on \(H\).

The increasing local spaces exhaust \(H\).
Because \(C\) is Hilbert--Schmidt, finite-rank approximation gives
\[
 \|C_L-C\|_{\mathcal S_2}\to0,\qquad
 \|M_LC_L\|_{\mathcal S_2}
       +\|C_LM_L\|_{\mathcal S_2}\to0 .
 \tag{PE.22}
\]
For the second statement, first use finite-rank approximation of
\(C\) and strong convergence of the uniformly bounded \(M_L\);
then absorb \(C_L-C\) in Hilbert--Schmidt norm.
The adjoint gives the other product.
This uses no relative second-eigenvalue estimate.

\proofstep{6. Only the odd relative gap enters the comparison.}
Use global parity spaces, extend \(M_{L,\pm}\) by zero and
\(S_\pm\) by the identity outside \(H_L^\pm\).
For \(0\le x\le1-\kappa_L\),
\[
 (1-x)^{-1/2}-1
       =\frac{x}{\sqrt{1-x}(1+\sqrt{1-x})},
 \qquad
 |\sqrt{1-x}-1|\le x .
\]
Functional calculus and (PE.7) therefore yield
\[
 \begin{aligned}
 \|B_L-C_L\|_{\mathcal S_2}
 \le\sqrt2\bigl(
      \kappa_L^{-1/2}\|M_{L,-}X_L\|_{\mathcal S_2}
          +\|X_LM_{L,+}\|_{\mathcal S_2}\bigr).
 \end{aligned}
 \tag{PE.23}
\]
For the second term use
\(S_+^{1/2}-I=-M_{L,+}(I+S_+^{1/2})^{-1}\);
the inverse here has norm at most one even on the kernel.
For the first term factor
\(S_-^{-1/2}X_LS_+^{1/2}-X_L\) as
\((S_-^{-1/2}-I)X_LS_+^{1/2}
 +X_L(S_+^{1/2}-I)\).
Thus no inverse even-gap bound is hidden in (PE.23).
Under (PE.8), the right side tends to zero by (PE.22).
Together with (PE.22) this proves the first part of (PE.9).

More generally, the same proof works if
\(\kappa_L^{-1/2}\|M_{L,-}X_L\|_{\mathcal S_2}\to0\).
This is the explicit comparison estimate, not a claim that the
displayed decay has been proved when \(\kappa_L\to0\).

\proofstep{7. Entire products with uniform tails.}
Let the positive eigenvalues of \(B_L\) be
\(\beta_{L,1}\ge\beta_{L,2}\ge\cdots>0\).
Reflection pairs positive and negative eigenvalues, and Step 4
makes every nonzero eigenvalue simple.
Hilbert--Schmidt convergence implies operator-norm convergence,
so the min--max principle gives, for every fixed \(j\),
\[
                    \beta_{L,j}\to1/\gamma_j .
 \tag{PE.24}
\]
Since the norms also converge and the eigenvalues are paired,
\[
  \sum_j\beta_{L,j}^2=\tfrac12\|B_L\|_{\mathcal S_2}^2
                   \longrightarrow
                    \sum_j\gamma_j^{-2}.
 \tag{PE.25}
\]
Subtracting the first \(J\) terms using (PE.24) proves
\[
 \limsup_{L\to\infty}\sum_{j>J}\beta_{L,j}^2
                         =\sum_{j>J}\gamma_j^{-2}\to0 .
 \tag{PE.26}
\]
No interlacing between operators in different metrics is used.

The ground Fourier transform is entire of order at most one,
even, and nonzero at zero.  Step 4 identifies all its zeros,
including their analytic multiplicities, as
\(\pm\beta_{L,j}^{-1}\).
Hadamard factorization and (PE.25) give
\[
 \frac{F(z)}{F(0)}
                =\prod_{j\ge1}(1-z^2\beta_{L,j}^2).
 \tag{PE.27}
\]
Pairing opposite zeros cancels genus-one exponentials;
evenness removes a remaining linear exponential, and the value
at zero removes the constant.
For \(|z|\le R\), choose \(J\) so that
\(\gamma_{J+1}>2R\).
Then (PE.24) makes the tail factors uniformly within the
convergent logarithm disk for all sufficiently large \(L\), and
\[
 \left|\sum_{j>J}\log(1-z^2\beta_{L,j}^2)\right|
                   \le 2R^2\sum_{j>J}\beta_{L,j}^2 .
\]
Equations (PE.24)--(PE.26) give local uniform convergence of
(PE.27) to the product over distinct positive ordinates.
If those zeros are all simple, the ordinary even Hadamard
factorization of \(\Xi\) identifies it with \(\Xi/\Xi(0)\).
This completes the proof.
\end{proof}

\paragraph{Exact comparison with the previous selection criterion.}
The old AC-S48 ratio \(q(v_\lambda)/a_2(\log\lambda)\to0\),
under \(A\), gives an eventual simple even minimum and
\(a_e/a_o\le a_1/a_2\to0\).  AC-S49 separately shows that
this same ratio forces simplicity of the zeta zeros.
Thus the old criterion implies the parity separation and
simplicity hypotheses used here.

The new proof does not supply either hypothesis from \(A\).
Parity separation alone yields \(\Delta_\Gamma\), not the
multiplicity-sensitive target if multiple zeros exist.
Adding explicit simplicity yields the original normalized
ground Fourier limit with only a fixed relative odd gap,
instead of the old prolate-trial ratio.
This is the precise conditional comparison; replacing the
target by its reduced zero product is not counted as closing
the original \(\Xi\) question.  Also, entire convergence in
(PE.9) does not by itself preclude escape of spatial \(L^2\)
mass, so the weighted and spatially unit-normalized conclusions
of AC-S47--AC-S48 are not silently imported.

\paragraph{Source and geometric scope.}
The proof uses the actual local operator/form and compact
resolvent of AC-S03 and the global sampling completion proved
directly in AC-S50.  Its primitive, rank-one commutator,
quotient selfadjointness, eigenvalue/evaluation correspondence,
parity-block comparison and uniform product tails are proved
above.  Suzuki's primary Sections 2--4 supply a parallel
unshifted inverse-compression description and were freshly read;
no equation here identifies the shifted quotient by analogy alone.
The relation to the literal finite S25 determinant requires its
already audited fixed-window scalar/phase normalization; no
new joint Galerkin/window rate is asserted.
These are Hilbert-space and entire-function statements, not
new etale cohomology, a motivic realization, or an arithmetic
square.
